% Generated draft: requires mathematical and source review.
\documentclass[11pt]{article}
\usepackage[margin=25mm]{geometry}
\usepackage{fontspec}
\setmainfont{Latin Modern Roman}
\usepackage{amsmath,amssymb,mathtools,unicode-math}
\setmathfont{Latin Modern Math}
\usepackage{enumitem,needspace,xurl,hyperref}
\setlength{\parindent}{0pt}
\setlength{\parskip}{5pt}
\setlength{\emergencystretch}{3em}
\newcommand{\N}{\mathbb{N}}
\newcommand{\Z}{\mathbb{Z}}
\newcommand{\Q}{\mathbb{Q}}
\newcommand{\R}{\mathbb{R}}
\newcommand{\C}{\mathbb{C}}
\newcommand{\F}{\mathbb{F}}
\newcommand{\A}{\mathbb{A}}
\newcommand{\PP}{\mathbb{P}}
\newcommand{\E}{\mathbb{E}}
\newcommand{\eps}{\varepsilon}
\newcommand{\dd}{\,\mathrm{d}}
\begin{document}
\section*{UnsolvedMath ID: 30000247}
\textbf{Source problem code:} catalogue entry
\subsection{Definitions and mathematical statement}
Let \(h\) be a multiplicative generalized cohomology theory and let \(r:h\to H\R\) be the given natural multiplicative transformation. On smooth manifolds \(X\), one standard meaning of a smooth extension is a differential refinement \(\widehat h^q(X)\) with natural maps
\[
I:\widehat h^q(X)\to h^q(X),\qquad
R:\widehat h^q(X)\to\Omega^q_{\mathrm{cl}}(X),\qquad
a:\Omega^{q-1}(X)/\operatorname{im}\dd\to\widehat h^q(X),
\]
satisfying \([R(\widehat x)]_{\mathrm{dR}}=r(I(\widehat x))\), \(R\circ a=\dd\), and the exact sequence
\[
h^{q-1}(X)\xrightarrow{r}\Omega^{q-1}(X)/\operatorname{im}\dd
\xrightarrow{a}\widehat h^q(X)\xrightarrow{I}h^q(X)\longrightarrow0.
\]
In the first arrow, a real cohomology class is viewed through the de Rham identification as a closed form modulo exact forms. These are conventional differential-cohomology axioms, not definitions reproduced in the supplied record. If \(H\R\) is meant with a graded coefficient algebra rather than scalar real coefficients, the differential-form complex must be changed accordingly.

The two source questions are to determine hypotheses on \((h,r)\) under which such a smooth extension exists, and conditions under which it is unique. A standard uniqueness interpretation is uniqueness up to natural isomorphism commuting with \(I,R,a\), rather than literal equality or uniqueness of the isomorphism itself. If multiplicativity of the extension is required, compatible products and preservation of those products must also be included. The given transformation is multiplicative, but the extracted text does not specify the full structure required of its smooth extension. Circle integration, a spectrum-level differential datum, and further coherence conditions are not imposed silently; including them can change uniqueness questions.

The supplied literature assessment checked 2026-08-21 reports existence in general frameworks and uniqueness for important theories including \(K\)-theory and bordism, but says the unrestricted extracted formulation is broader than those results. The exact hypotheses and notion of uniqueness in those theorems are not supplied, so no unrestricted uniqueness theorem is asserted.
\subsection{Short English statement}
Determine when a multiplicative cohomology transformation to real cohomology admits a differential refinement and when that refinement is unique in the intended category. The precise smooth-extension axioms and extra structures must be fixed before applying existence or uniqueness results.
\subsection{Sources}
\begin{itemize}
\item Local UnsolvedMath registry, permanent ID 30000247, catalogue entry, \emph{Smooth Extensions of Multiplicative Cohomology Transformations}.
\item \emph{Geometric Topology and Connections with Quantum Field Theory} (2006), Oberwolfach Reports, pp. 1505--1546, DOI \url{https://doi.org/10.4171/owr/2005/27}.
\item Supplied reference labeled ``Bunke uniqueness theorem'': \url{https://doi.org/10.1112/jtopol/jtq002}.
\item Supplied reference labeled ``Smooth K-theory survey'': \url{https://citeseerx.ist.psu.edu/document?doi=e1cab8a8f9f5ce8e734b609e49f6f08629ae7a3&repid=rep1&type=pdf}. The query ampersands are rendered as URL characters from the source's escaped form.
\item Supplied literature triage checked 2026-08-21; theorem hypotheses and full secondary bibliographic entries are not reproduced.
\end{itemize}
\end{document}
